\section{Conclusion}
\label{sec:conclusion}

Serving agents is not serving requests. A session arrives with state, leaves it behind, and
comes back to it, and on a single GPU the aggregate of that state exceeds the pool by an
order of magnitude. The systems we have are built around a threshold---how many, how much
concurrency, how high a watermark---which forces the operator to guess at a future the
workload does not reveal until it happens.

Once a session can be suspended and revived, the guess is unnecessary. Capacity becomes a
ledger of bytes, admission becomes a feasibility test against that ledger, and the two
decisions that remain---whose state deserves a slot, and whether to copy or recompute---are
priced in measured time. The design needs no predictor of the next round, and the numbers
say why: the part of the footprint that would have to be predicted is a few percent of the
part that is measured.

We have instantiated the principle as a controller outside an unmodified serving engine and
measured it against request-level serving and an agent-level feedback controller on a
19-session coding-agent replay. \pend{one-sentence result with the three arms}. The
reusable part is not the constants but the shape of the argument: state the feasible set,
price the recovery, and let the workload decide---rather than tune a threshold and hope the
next round looks like the last.
